Physical land use---the actual pattern of residential neighborhoods,
commercial corridors, industrial zones, open land, and water bodies on
the ground---is among the oldest and most intuitive bases for community
of interest recognition in American redistricting.
Yet algorithmic redistricting systems have historically operationalized
geographic community only through boundary length, ignoring the land cover
character of the tracts on either side of each boundary.
This paper presents M.2, the second paper in the M-track community character
series, which operationalizes physical land use through the USGS National
Land Cover Database (NLCD) 2021 30-meter raster.

For each census tract we compute a four-fraction land use character vector:
fraction of pixels classified as residential (NLCD classes 22--23), commercial
and industrial (class 24), open developed (class 21), and natural (all other
non-water classes including forest, wetland, shrub, and agricultural).
Edge weights between adjacent tracts are modulated by the cosine similarity
of their land use vectors: residentially similar tracts receive heavier edges;
residential-to-commercial transitions receive lighter edges that serve as
natural district cut seams.
A hard-cut rule enforces water boundaries: any adjacency where both tracts
have majority water classification (class 11) receives an edge weight of zero,
preventing the algorithm from assigning cross-water tracts to the same district
in the absence of explicit bridge adjacency.

Experimental design targets North Carolina ($k=14$), Wisconsin ($k=8$), and
Texas ($k=38$) with the primary metric of fraction of district boundaries
coinciding with NLCD land-use class transitions.
Compactness comparison (Polsby-Popper) is the secondary metric.
All quantitative results are deferred to Phase~2 pending NLCD GeoTIFF
download and zonal statistics computation.
The NLCD data is annual, free, publicly available, and contains no
racial or partisan content, making physical land use one of the most
legally defensible community-of-interest signals available.
